\scanpage{114}
alors $\exists$ \inserted{local\textsuperscript{t}} autom \deleted{\unsure{int.}}\inserted{$\theta$} de $A$, tel que
\[
\theta(H)=\overline H)\qquad\theta(H')=\overline{H'}
\]
ssi
\[
\overline{\rho_0}=\rho_0
\]

D'autre part, si on se donne $\rho_0$ \inserted{(inv)},  il existe \uline{\unsure{loc|alors|local\textsuperscript{t}}} une paire $(H,H')$ qui le réalise.

D'ailleurs, \deleted{tout autom de $A$ \deleted{\unsure{qui}} \inserted{est} \deleted{\illegible{2}} intérieur loc. Zar (et même glob.\textsuperscript{ment}, si $\operatorname{Pic}(k)=0$), et s'il fixe $H$ il \deleted{est loc. \unsure{un|de}} \inserted{provient loc. Zar} d'une section de $H$ ($H$ est son propre centralisateur dans $A$) \deleted{\unsure{ou} est} provient d'un $\theta\in$ \deleted{de la \unsure{façon|forme}} $x\longmapsto x'=\operatorname{Tr}(x)1-x$}

\deleted{La \unsure{conjug.} par un tel autom provient de} Dans $G=\underline{\operatorname{Aut}}(A)\simeq\underline{\operatorname{Aut}}(P)$, le sous-schéma en groupes qui stabilise à la fois $H$ et $H'$ est un sous-\deleted{gr}\inserted{schéma}
% slanted
\margin{du moins si $\rho\ne-2$ \deleted{partout} i.e. $\rho_0\ne-4$ \deleted{partout}, sinon il y a un groupe $\unsure{C|\overline C}$ d'autos qui est $\simeq\mathbb Z/2\mathbb Z\times\mathbb Z/2\mathbb Z$.} en groupes isom. à $(\mathbb Z/2\mathbb Z)_S$, \deleted{\unsure{qui}} dont la section non nulle est un auto- qui, sur $H$ et sur $H'$ \deleted{est donné par}
% archive sheet number: 113
